\documentclass[UTF8,12pt,a4paper,fontset=fandol]{ctexart}

% 支持从项目根目录或本文件所在目录编译。
\makeatletter
\def\input@path{{../}{../../}}
\makeatother
\usepackage{researchnotes}

\title{行列式自测卷（含答案）}
\author{}
\date{}

\begin{document}
\maketitle

\section*{作答说明}
本卷面向高等代数行列式章节的阶段复习，共 30 题，满分 100 分。
第一部分每题 2 分，第二部分每题 3 分，第三部分每题 5 分。
建议用时 180--240 分钟，也可分三次完成。计算题应写出主要变换，证明题应交代关键依据。
除题目另有说明外，所有矩阵及参数均在实数域 $\mathbb R$ 上，$n$ 为正整数。

记 $I_n$ 为 $n$ 阶单位矩阵，$A^{\mathsf T}$ 为 $A$ 的转置。
对方阵 $A=(a_{ij})$，记 $M_{ij}$ 为删去第 $i$ 行、第 $j$ 列所得的
\keyterm{余子式}（minor），$C_{ij}=(-1)^{i+j}M_{ij}$ 为对应的
\keyterm{代数余子式}（cofactor）。$R_i$、$K_j$ 分别表示第 $i$ 行、第 $j$ 列。
空乘积约定为 $1$，零次幂表示常数多项式 $1$ 的取值。
可使用行列式基本性质、乘法公式、拉普拉斯展开和范德蒙公式；最后两题还需用到一元多项式的因式定理。

% ANSWER-INTRO-BEGIN
各题后附参考解答，等价的正确方法同样有效。
自测时建议先使用题目版独立作答，再核对本文件中的推导。
% ANSWER-INTRO-END

\section{基础概念与性质（共 20 分）}
\begin{enumerate}[label=\arabic*.,series=detquestions]

\item\label{q:inversions}
求排列 $(3,6,1,5,2,4)$ 的\keyterm{逆序数}（inversion number），并判断其奇偶性。
\begin{solve}
逐个数出每个元素右侧比它小的元素个数，依次为 $2,4,0,2,0,0$。
因此逆序数为 $8$，这是偶排列，其符号为 $(-1)^8=1$。
\end{solve}

\item\label{q:term-sign}
在一般五阶行列式 $\det(a_{ij})_{1\leq i,j\leq5}$ 的完全展开式中，
单项式 $a_{14}a_{21}a_{35}a_{42}a_{53}$ 的系数是多少？说明依据。
\begin{solve}
行指标已按 $1,2,3,4,5$ 排列，列指标排列为 $(4,1,5,2,3)$。
其逆序数为 $3+0+2+0+0=5$，故该项的系数为 $(-1)^5=-1$。
\end{solve}

\item\label{q:count-terms}
将一般 $n$ 阶行列式看作关于彼此独立的元素 $a_{ij}$ 的多项式，设 $n\geq2$。
完全展开式中含 $a_{12}$ 的项有多少个？同时含 $a_{12}$ 和 $a_{21}$ 的项有多少个？
\begin{solve}
固定 $a_{12}$ 后，其余 $n-1$ 行与其余 $n-1$ 列可任意一一配对，故有 $(n-1)!$ 项。
再固定 $a_{21}$ 后，剩余配对数为 $(n-2)!$，其中 $n=2$ 时 $0!=1$。
\end{solve}

\item\label{q:row-operations}
设三阶矩阵 $A$ 满足 $\det A=5$。从 $A$ 出发，依次交换第一、三行，
将第二行乘以 $-2$，再把第一行的 $3$ 倍加到第三行，所得矩阵记为 $B$。求 $\det B$。
\begin{solve}
交换两行使行列式变号，一行乘以 $-2$ 使行列式乘以 $-2$，
最后的倍加变换不改变行列式。因此
\[
\det B=(-1)(-2)\det A=10.
\]
\end{solve}

\item\label{q:minor-cofactor}
设
\[
A=\begin{pmatrix}1&2&0\\3&-1&4\\2&5&1\end{pmatrix}.
\]
求 $M_{23}$ 和 $C_{23}$，并写出二者的符号关系。
\begin{solve}
删去第二行、第三列后，有
\[
M_{23}=\begin{vmatrix}1&2\\2&5\end{vmatrix}=1,
\qquad C_{23}=(-1)^{2+3}M_{23}=-1.
\]
余子式本身不包含位置符号 $(-1)^{i+j}$，代数余子式包含该符号。
\end{solve}

\item\label{q:cofactor-sums}
设三阶矩阵 $A=(a_{ij})$ 满足 $\det A=7$。求
\[
\sum_{j=1}^{3}a_{2j}C_{2j},\qquad
\sum_{j=1}^{3}a_{1j}C_{2j},
\]
并简述理由。
\begin{solve}
第一个和是原行列式按第二行展开，等于 $7$。
第二个和是把第二行换成第一行所得行列式按第二行的展开式：
替换第二行不会改变 $C_{2j}$，而新矩阵有两行相同，故该和为 $0$。
\end{solve}

\item\label{q:triangular}
计算下列两个行列式，并说明第二个行列式的符号：
\[
D_1=\begin{vmatrix}-2&1&0&3\\0&3&4&1\\0&0&1&5\\0&0&0&4\end{vmatrix},
\qquad
D_2=\begin{vmatrix}0&0&0&1\\0&0&2&0\\0&3&0&0\\4&0&0&0\end{vmatrix}.
\]
\begin{solve}
上三角行列式等于主对角线元素之积，故 $D_1=(-2)\cdot3\cdot1\cdot4=-24$。
第二个行列式仅有反序排列 $(4,3,2,1)$ 对应的项可能非零，其逆序数为 $6$，故
\[
D_2=(-1)^6(1\cdot2\cdot3\cdot4)=24.
\]
一般 $n$ 阶反对角行列式的符号因子为 $(-1)^{n(n-1)/2}$。
\end{solve}

\item\label{q:dependent-rows}
设 $D_n=\det(i+j)_{1\leq i,j\leq n}$，$n\geq2$。分别求 $n=2$ 与 $n\geq3$ 时的 $D_n$。
\begin{solve}
当 $n=2$ 时，$D_2=2\cdot4-3\cdot3=-1$。
当 $n\geq3$ 时，对每个列指标 $j$ 都有
\[
(1+j)-2(2+j)+(3+j)=0,
\]
即 $R_1-2R_2+R_3=0$。用两次倍加变换将第三行化成零行，得 $D_n=0$。
\end{solve}

\item\label{q:true-false}
判断下列断言的正误；正确的说明理由，错误的给出反例。
\begin{enumerate}[label=(\arabic*)]
\item 对任意同阶方阵 $A,B$，都有 $\det(A+B)=\det A+\det B$。
\item 一个三阶行列式等于零，则它一定有两行成比例。
\item 对任意实方阵 $A$，都有 $\det(A^{\mathsf T}A)\geq0$。
\end{enumerate}
\begin{solve}
(1) 错误。取 $A=B=I_2$，则 $\det(A+B)=4$，而 $\det A+\det B=2$。
行列式只在其他行（或列）固定时，对某一行（或列）具有线性性质。

(2) 错误。例如
\[
A=\begin{pmatrix}1&0&1\\0&1&1\\1&1&2\end{pmatrix}
\]
的第三行是前两行之和，故行列式为零，但任意两行均不成比例。

(3) 正确。由乘法公式与转置性质，
$\det(A^{\mathsf T}A)=\det(A^{\mathsf T})\det A=(\det A)^2\geq0$。
这里使用了 $\det A$ 为实数。
\end{solve}

\item\label{q:product-scaling}
设 $A,B$ 均为三阶矩阵，$\det A=-2$，$\det B=3$。求
\[
\det(2A),\qquad \det(A^{\mathsf T}B),\qquad \det(A^2).
\]
\begin{solve}
对三阶矩阵整体乘以 $2$，相当于三行各乘以 $2$。因此
\[
\det(2A)=2^3(-2)=-16,\qquad
\det(A^{\mathsf T}B)=(-2)\cdot3=-6,\qquad
\det(A^2)=(-2)^2=4.
\]
\end{solve}
\end{enumerate}

\section{计算与方法训练（共 30 分）}
\begin{enumerate}[label=\arabic*.,resume=detquestions]

\item\label{q:numeric-three}
用初等行变换计算
\[
D=\begin{vmatrix}1&2&3\\2&5&7\\1&0&4\end{vmatrix}.
\]
\begin{solve}
依次作 $R_2\leftarrow R_2-2R_1$、$R_3\leftarrow R_3-R_1$，得
\[
D=\begin{vmatrix}1&2&3\\0&1&1\\0&-2&1\end{vmatrix}
=\begin{vmatrix}1&1\\-2&1\end{vmatrix}=3.
\]
\end{solve}

\item\label{q:numeric-four}
计算
\[
D=\begin{vmatrix}
1&1&1&1\\1&2&1&1\\1&1&3&1\\1&1&1&4
\end{vmatrix}.
\]
\begin{solve}
将第二、三、四行各减去第一行，得到
\[
D=\begin{vmatrix}
1&1&1&1\\0&1&0&0\\0&0&2&0\\0&0&0&3
\end{vmatrix}=6.
\]
\end{solve}

\item\label{q:parameter-three}
将下列行列式分解为一次因式的乘积，并求它为零时 $x$ 的全部实数取值：
\[
D(x)=\begin{vmatrix}x&1&1\\1&x&1\\1&1&x\end{vmatrix}.
\]
\begin{solve}
作 $K_1\leftarrow K_1+K_2+K_3$，再从第一列提取公因子 $x+2$，有
\[
D(x)=(x+2)\begin{vmatrix}1&1&1\\1&x&1\\1&1&x\end{vmatrix}.
\]
将第二、三行分别减去第一行，得
\[
D(x)=(x+2)(x-1)^2.
\]
故零点为 $x=1$ 或 $x=-2$。提取公因子使用的是线性恒等式，
没有除以 $x+2$，因此在 $x=-2$ 时仍成立。
\end{solve}

\item\label{q:vandermonde-numeric}
利用\keyterm{范德蒙行列式}（Vandermonde determinant）计算
\[
D=\begin{vmatrix}
1&1&1&1\\-1&0&2&3\\1&0&4&9\\-1&0&8&27
\end{vmatrix}.
\]
\begin{solve}
四列对应节点 $-1,0,2,3$，各行的幂次依次为 $0,1,2,3$。
按“后列节点减前列节点”的方向，有
\[
D=(0+1)(2+1)(3+1)(2-0)(3-0)(3-2)=72.
\]
\end{solve}

\item\label{q:polynomial-four}
设 $x_1,x_2,x_3,x_4\in\mathbb R$。将下列行列式写成乘积形式：
\[
D=\begin{vmatrix}
1&1&1&1\\
x_1+1&x_2+1&x_3+1&x_4+1\\
x_1^2-x_1+2&x_2^2-x_2+2&x_3^2-x_3+2&x_4^2-x_4+2\\
x_1^3+2x_1&x_2^3+2x_2&x_3^3+2x_3&x_4^3+2x_4
\end{vmatrix}.
\]
\begin{solve}
先作 $R_2\leftarrow R_2-R_1$，把第二行化为 $(x_1,x_2,x_3,x_4)$。
再使用更新后的第二行，作 $R_3\leftarrow R_3+R_2-2R_1$、
$R_4\leftarrow R_4-2R_2$，得到标准范德蒙行列式。因此
\[
D=\prod_{1\leq i<j\leq4}(x_j-x_i).
\]
这些变换不要求节点互异；有节点重复时，等式两边均为零。
\end{solve}

\item\label{q:constant-off-diagonal}
设 $n\geq2$，$a,b\in\mathbb R$，矩阵 $A_n=(a_{ij})$ 的元素为
\[
a_{ij}=\begin{cases}a,&i=j,\\b,&i\ne j.\end{cases}
\]
求 $\det A_n$，并说明所得公式在 $a=b$ 时是否成立。
\begin{solve}
将其余各列加到第一列，第一列各元素均为 $s=a+(n-1)b$。
从第一列提取 $s$ 后，该列全为 $1$；再将第 $2$ 至第 $n$ 行各减去第一行。
所得矩阵的第一列为 $(1,0,\ldots,0)^{\mathsf T}$，右下角为 $(a-b)I_{n-1}$，故
\[
\boxed{\det A_n=\bigl(a+(n-1)b\bigr)(a-b)^{n-1}.}
\]
推导中没有除以 $s$ 或 $a-b$。当 $a=b$ 时各行相同，行列式为零，公式仍然成立。
\end{solve}

\item\label{q:arrow}
设 $n\geq2$，计算下列 $n$ 阶行列式，其中 $d_2,\ldots,d_n$ 允许为零：
\[
D_n=\begin{vmatrix}
\alpha&\beta_2&\beta_3&\cdots&\beta_n\\
\gamma_2&d_2&0&\cdots&0\\
\gamma_3&0&d_3&\cdots&0\\
\vdots&\vdots&\vdots&\ddots&\vdots\\
\gamma_n&0&0&\cdots&d_n
\end{vmatrix}.
\]
\begin{solve}
使用排列展开。若第一行取第一列，则其他行必须取各自的对角元，贡献为
$\alpha\prod_{i=2}^n d_i$。
若第一行取第 $k$ 列（$k\geq2$），则第 $k$ 行只能取第一列，其他行取对角元。
对应排列是换位 $(1\ k)$，符号为负。因此
\[
\boxed{D_n=\alpha\prod_{i=2}^n d_i
-\sum_{k=2}^n\beta_k\gamma_k
\prod_{\substack{2\leq i\leq n\\i\ne k}}d_i.}
\]
上述两类已包含所有可能非零的排列项；推导没有作除法，所以允许任意 $d_i=0$。
当 $n=2$ 时，求和项内的空乘积为 $1$。
\end{solve}

\item\label{q:tridiagonal}
设 $T_n=(t_{ij})_{1\leq i,j\leq n}$，其中
\[
t_{ij}=\begin{cases}2,&i=j,\\-1,&|i-j|=1,\\0,&|i-j|\geq2.\end{cases}
\]
记 $D_n=\det T_n$。建立递推关系并求 $D_n$ 的通项。
\begin{solve}
为统一递推，约定 $D_0=1$，并有 $D_1=2$。
当 $n\geq2$ 时按第一行展开。第一项为 $2D_{n-1}$；
第二项中，元素 $-1$ 与位置符号 $-1$ 相乘得 $1$，
其余子式的第一列仅有首项 $-1$，再展开得 $-D_{n-2}$。因此
\[
D_n=2D_{n-1}-D_{n-2}\qquad(n\geq2).
\]
于是 $D_n-D_{n-1}=D_{n-1}-D_{n-2}$，各相邻差均为 $D_1-D_0=1$。
从 $D_0=1$ 出发累加，得到
\[
\boxed{D_n=n+1.}
\]
\end{solve}

\item\label{q:cyclic}
设 $n\geq3$，矩阵 $A=(a_{ij})$ 仅在以下位置可能非零：
\[
a_{ii}=u_i\ (1\leq i\leq n),\qquad
a_{i,i+1}=v_i\ (1\leq i<n),\qquad a_{n1}=v_n.
\]
其余元素均为零。求 $\det A$。
\begin{solve}
若排列项全部取对角元，其贡献为 $\prod_{i=1}^n u_i$。
若有一行取到后一列的 $v_i$（列指标按模 $n$ 循环），
该列被占用就迫使下一行也取后一列，依次迫使所有行都取 $v_i$。
因此仅有两种排列可能产生非零项。第二种的列指标排列为 $(2,3,\ldots,n,1)$，
逆序数为 $n-1$，所以
\[
\boxed{\det A=\prod_{i=1}^n u_i+(-1)^{n-1}\prod_{i=1}^n v_i.}
\]
即使某些 $u_i$ 或 $v_i$ 为零，这一排列展开仍成立。
\end{solve}

\item\label{q:laplace}
使用\keyterm{拉普拉斯展开}（Laplace expansion），按第一、三行展开并计算
\[
D=\begin{vmatrix}
1&2&0&0\\0&0&3&4\\5&6&0&0\\0&0&7&8
\end{vmatrix}.
\]
请写出所选子式及其代数余子式的符号。
\begin{solve}
在第一、三行中选两列组成二阶子式，只有列集合 $\{1,2\}$ 可能产生非零子式。
其补余子式由第二、四行及第三、四列组成。拉普拉斯展开的符号为
$(-1)^{(1+3)+(1+2)}=-1$，因此
\[
D=-\begin{vmatrix}1&2\\5&6\end{vmatrix}
\begin{vmatrix}3&4\\7&8\end{vmatrix}
=-(-4)(-4)=-16.
\]
所选二阶子式为 $-4$，它的代数余子式为 $-(-4)=4$。
\end{solve}
\end{enumerate}

\section{综合应用与证明（共 50 分）}
\begin{enumerate}[label=\arabic*.,resume=detquestions]

\item\label{q:cramer}
用\keyterm{克拉默法则}（Cramer's rule）求解方程组，写出系数行列式及三个替换列行列式：
\[
\begin{cases}
x+y+z=6,\\
2x-y+3z=9,\\
3x+2y-z=4.
\end{cases}
\]
\begin{solve}
系数行列式为
\[
D=\begin{vmatrix}1&1&1\\2&-1&3\\3&2&-1\end{vmatrix}
=(-5)-(-11)+7=13\ne0.
\]
分别用常数列替换第一、二、三列，得到
\begin{align*}
D_x&=\begin{vmatrix}6&1&1\\9&-1&3\\4&2&-1\end{vmatrix}
=-30+21+22=13,\\
D_y&=\begin{vmatrix}1&6&1\\2&9&3\\3&4&-1\end{vmatrix}
=-21+66-19=26,\\
D_z&=\begin{vmatrix}1&1&6\\2&-1&9\\3&2&4\end{vmatrix}
=-22+19+42=39.
\end{align*}
由克拉默法则，方程组有唯一解
\[
\boxed{(x,y,z)=(1,2,3).}
\]
代回三式，左端依次为 $6,9,4$，与常数项一致。
\end{solve}

\item\label{q:parameter-system}
讨论实参数 $\lambda$ 取不同值时，下列方程组的解的情况；有解时写出全部解：
\[
\begin{cases}
\lambda x+y+z=1,\\
x+\lambda y+z=1,\\
x+y+\lambda z=-2.
\end{cases}
\]
\begin{solve}
由第~\ref{q:parameter-three} 题，系数行列式为
$D=(\lambda-1)^2(\lambda+2)$。

当 $\lambda\notin\{1,-2\}$ 时，$D\ne0$，所以有唯一解。
令 $s=x+y+z$，三式可写为
\[
(\lambda-1)x+s=1,\quad
(\lambda-1)y+s=1,\quad
(\lambda-1)z+s=-2.
\]
三式相加得 $(\lambda+2)s=0$，故 $s=0$，从而
\[
\boxed{(x,y,z)=\left(\frac1{\lambda-1},\frac1{\lambda-1},
-\frac2{\lambda-1}\right).}
\]

当 $\lambda=1$ 时，前两式要求 $x+y+z=1$，第三式要求 $x+y+z=-2$，故无解。

当 $\lambda=-2$ 时，三式分别给出
$x=(s-1)/3$、$y=(s-1)/3$、$z=(s+2)/3$，它们的和恒为 $s$。
因此 $s$ 可以任取实数，全部解为
\[
\boxed{(x,y,z)=(-1/3,-1/3,2/3)+t(1,1,1),\qquad t\in\mathbb R.}
\]
两个使系数行列式为零的参数对应不同情况，不能仅凭 $D=0$ 判定无解。
\end{solve}

\item\label{q:all-cofactors}
设 $n\geq2$，$J_n$ 为所有元素都等于 $1$ 的 $n$ 阶矩阵。
对任意 $n$ 阶矩阵 $A$，证明
\[
\det(A+tJ_n)=\det A+t\sum_{i=1}^n\sum_{j=1}^n C_{ij}
\qquad(t\in\mathbb R),
\]
其中 $C_{ij}$ 始终指原矩阵 $A$ 的代数余子式。进而求
\[
A=\begin{pmatrix}2&1&0\\0&3&1\\1&0&2\end{pmatrix}
\]
的全部代数余子式之和，要求不逐个计算九个代数余子式。
\begin{solve}
记 $A$ 的各列为 $a_1,\ldots,a_n$，$e=(1,\ldots,1)^{\mathsf T}$。
按列的多重线性性质展开 $\det(a_1+te,\ldots,a_n+te)$。
凡从至少两列中选取 $te$ 的项，都因两列相同而为零。
没有选取 $te$ 的项为 $\det A$；仅在第 $j$ 列选取 $te$ 的项为
\[
t\det(a_1,\ldots,a_{j-1},e,a_{j+1},\ldots,a_n)
=t\sum_{i=1}^n C_{ij}.
\]
相加即得所证恒等式，它不要求 $A$ 可逆。

对题给矩阵，在 $A+tJ_3$ 中作 $R_2\leftarrow R_2-R_1$、
$R_3\leftarrow R_3-R_1$，得到
\[
\det(A+tJ_3)=
\begin{vmatrix}2+t&1+t&t\\-2&2&1\\-1&-1&2\end{vmatrix}
=5(2+t)+3(1+t)+4t=13+12t.
\]
比较 $t$ 的系数，全部代数余子式之和为 $\boxed{12}$。
\end{solve}

\item\label{q:skew-symmetric}
设实方阵 $A$ 满足 $A^{\mathsf T}=-A$，称其为
\keyterm{反对称矩阵}（skew-symmetric matrix）。
证明：若 $A$ 的阶数为奇数，则 $\det A=0$。
再给出一个偶数阶反对称矩阵，其行列式不为零。
\begin{solve}
若 $A$ 为 $n$ 阶矩阵，则
\[
\det A=\det(A^{\mathsf T})=\det(-A)=(-1)^n\det A.
\]
当 $n$ 为奇数时，得 $2\det A=0$。由于在实数域上讨论，故 $\det A=0$。
偶数阶反例可取
\[
A=\begin{pmatrix}0&1\\-1&0\end{pmatrix},
\qquad A^{\mathsf T}=-A,\qquad\det A=1.
\]
因此“阶数为奇数”是该结论的必要限制之一，不能删去。
\end{solve}

\item\label{q:block-determinant}
设 $A,B$ 为任意两个 $n$ 阶实矩阵。用行、列变换证明
\[
\det\begin{pmatrix}A&B\\B&A\end{pmatrix}
=\det(A+B)\det(A-B).
\]
不假设 $A$、$B$ 可逆，也不假设 $AB=BA$。
\begin{solve}
对每个 $i=1,\ldots,n$，将第 $n+i$ 行加到第 $i$ 行，得到
\[
\begin{pmatrix}A+B&A+B\\B&A\end{pmatrix}.
\]
再对每个 $j=1,\ldots,n$，将第 $n+j$ 列减去第 $j$ 列，得到
\[
\begin{pmatrix}A+B&0\\B&A-B\end{pmatrix}.
\]
这些都是倍加变换，不改变行列式。对最后的矩阵按前 $n$ 行作拉普拉斯展开，
只有前 $n$ 列对应的子式可能有贡献；位置符号为
$(-1)^{2(1+\cdots+n)}=1$。因此其行列式等于
$\det(A+B)\det(A-B)$，结论成立。
\end{solve}

\item\label{q:rank-one}
设 $u=(u_1,\ldots,u_n)^{\mathsf T}$、$v=(v_1,\ldots,v_n)^{\mathsf T}$ 为实列向量，
$uv^{\mathsf T}$ 表示第 $(i,j)$ 个元素为 $u_iv_j$ 的矩阵。
用行列式的多重线性性质证明
\[
\det(I_n+uv^{\mathsf T})=1+\sum_{i=1}^n u_iv_i.
\]
进而说明，对任意实向量 $u$，$\det(I_n+uu^{\mathsf T})>0$。
\begin{solve}
记 $e_j$ 为第 $j$ 个标准单位列向量。所求行列式为
\[
\det(e_1+v_1u,\ldots,e_n+v_nu).
\]
完全按列展开后，选取至少两列为 $v_ju$ 的项均因列成比例而为零。
不选取 $u$ 的项为 $1$；仅在第 $j$ 列选取 $v_ju$ 的项为
\[
v_j\det(e_1,\ldots,e_{j-1},u,e_{j+1},\ldots,e_n)=v_ju_j.
\]
最后一步可由 $u=\sum_i u_i e_i$ 按第 $j$ 列展开得到：
当 $i\ne j$ 时出现重复列，只有 $i=j$ 的项保留。
相加即得恒等式。取 $v=u$，有
\[
\det(I_n+uu^{\mathsf T})=1+\sum_{i=1}^n u_i^2\geq1>0.
\]
\end{solve}

\item\label{q:min-index}
设 $s_1,\ldots,s_n\in\mathbb R$，$A=(s_{\min(i,j)})_{1\leq i,j\leq n}$。
求 $\det A$，并给出它不为零的充要条件。
例如 $n=4$ 时，
\[
A=\begin{pmatrix}
s_1&s_1&s_1&s_1\\
s_1&s_2&s_2&s_2\\
s_1&s_2&s_3&s_3\\
s_1&s_2&s_3&s_4
\end{pmatrix}.
\]
\begin{solve}
按 $i=n,n-1,\ldots,2$ 的顺序作 $R_i\leftarrow R_i-R_{i-1}$。
由下往上操作，保证每次使用的第 $i-1$ 行仍为原行。
变换后，第 $i$ 行在 $j<i$ 时为零，在 $j\geq i$ 时为 $s_i-s_{i-1}$；
第一行保持不变。因此所得矩阵为上三角矩阵，故
\[
\boxed{\det A=s_1\prod_{i=2}^n(s_i-s_{i-1}).}
\]
从而 $\det A\ne0$ 当且仅当 $s_1\ne0$ 且每一对相邻参数都不同，
即 $s_i\ne s_{i-1}$（$2\leq i\leq n$）。
这里不要求所有参数两两不同。$n=1$ 时乘积为空，条件退化为 $s_1\ne0$。
\end{solve}

\item\label{q:polynomial-basis}
设 $p_k(t)=\sum_{r=0}^{k}c_{kr}t^r$（$k=0,1,\ldots,n-1$）为实系数多项式，
允许 $c_{kk}=0$。对任意 $x_1,\ldots,x_n\in\mathbb R$，证明
\[
\det\bigl(p_{i-1}(x_j)\bigr)_{1\leq i,j\leq n}
=\left(\prod_{k=0}^{n-1}c_{kk}\right)
\prod_{1\leq i<j\leq n}(x_j-x_i).
\]
进而求
\[
\begin{vmatrix}
2&2&2\\
3a-1&3b-1&3c-1\\
5a^2+a+4&5b^2+b+4&5c^2+c+4
\end{vmatrix}.
\]
\begin{solve}
令 $V=(x_j^{i-1})_{1\leq i,j\leq n}$，并令 $L$ 为下三角系数矩阵：
当 $0\leq r\leq k\leq n-1$ 时，$L_{k+1,r+1}=c_{kr}$，其余元素为零。
逐项计算矩阵乘积，得
\[
(LV)_{k+1,j}=\sum_{r=0}^{k}c_{kr}x_j^r=p_k(x_j).
\]
因此目标矩阵等于 $LV$。利用乘法公式和范德蒙公式，得到
\[
\det(LV)=\det L\det V
=\left(\prod_{k=0}^{n-1}c_{kk}\right)
\prod_{i<j}(x_j-x_i).
\]
这一论证没有对任何系数或节点差作除法，所以也覆盖最高次系数为零、节点重复等情形。
具体行列式对应的对角系数为 $2,3,5$，故其值为
\[
\boxed{30(b-a)(c-a)(c-b).}
\]
\end{solve}

\item\label{q:missing-power}
设 $n\geq2$，$x_1,\ldots,x_n\in\mathbb R$。
矩阵 $W$ 的前 $n-1$ 行依次为 $(x_j^k)_{j=1}^n$（$k=0,\ldots,n-2$），
最后一行为 $(x_j^n)_{j=1}^n$，即跳过了幂次 $n-1$。
证明
\[
\det W=\left(\sum_{j=1}^n x_j\right)
\prod_{1\leq i<j\leq n}(x_j-x_i).
\]
例如 $n=2$ 时，$W=\begin{pmatrix}1&1\\x_1^2&x_2^2\end{pmatrix}$。
\begin{solve}
记 $s=\sum_{j=1}^n x_j$，展开多项式
\[
P(t)=\prod_{j=1}^n(t-x_j)
=t^n-st^{n-1}+\sum_{k=0}^{n-2}b_kt^k.
\]
对每个 $j$，由 $P(x_j)=0$ 得
\[
x_j^n=sx_j^{n-1}-\sum_{k=0}^{n-2}b_kx_j^k.
\]
所以 $W$ 的最后一行等于 $s(x_1^{n-1},\ldots,x_n^{n-1})$
与前面各行的一个线性组合之和。
按最后一行的线性性质展开，来自前面各行的项均因出现重复行而为零，故
\[
\det W=s\det(x_j^{i-1})_{1\leq i,j\leq n}
=s\prod_{i<j}(x_j-x_i).
\]
上述多项式恒等式不要求节点互异，也不要求 $s\ne0$。
\end{solve}

\item\label{q:interpolation-sum}
设 $n\geq2$，$x_1,\ldots,x_n$ 是两两不同的实数。
证明：对每个整数 $m$，$0\leq m\leq n-1$，都有
\[
\sum_{k=1}^n
\frac{x_k^m}{\displaystyle\prod_{\substack{1\leq j\leq n\\j\ne k}}(x_k-x_j)}
=\begin{cases}0,&0\leq m\leq n-2,\\1,&m=n-1.\end{cases}
\]
要求用范德蒙行列式及代数余子式展开证明，并说明为什么需要节点两两不同。
\begin{solve}
记 $V=(x_j^{i-1})_{1\leq i,j\leq n}$，
$\Delta=\det V=\prod_{i<j}(x_j-x_i)\ne0$，并记 $C_{nk}$ 为 $V$ 的代数余子式。
将 $V$ 的最后一行换成 $(x_1^m,\ldots,x_n^m)$，所得行列式记为 $D_m$。
若 $m\leq n-2$，它与前面某行重复，故 $D_m=0$；若 $m=n-1$，则 $D_m=\Delta$。
替换最后一行不改变这一行对应的代数余子式，故
\[
D_m=\sum_{k=1}^n x_k^m C_{nk}.
\]
删去最后一行和第 $k$ 列后，余下的是 $n-1$ 阶范德蒙行列式，因而
\[
C_{nk}=(-1)^{n+k}
\prod_{\substack{1\leq i<j\leq n\\i,j\ne k}}(x_j-x_i).
\]
另一方面，$\Delta$ 中含节点 $x_k$ 的因子之积为
\[
\prod_{i<k}(x_k-x_i)\prod_{j>k}(x_j-x_k)
=(-1)^{n-k}\prod_{j\ne k}(x_k-x_j).
\]
两式相除时，符号相消，因为 $(n+k)-(n-k)=2k$ 为偶数。因此
\[
\frac{C_{nk}}{\Delta}=\frac1{\prod_{j\ne k}(x_k-x_j)}.
\]
将 $D_m$ 的展开式除以 $\Delta$，即得所证公式。
节点两两不同同时保证 $\Delta\ne0$ 和题目中所有分母非零；节点重复时原表达式不再有定义。
\end{solve}
\end{enumerate}

% ANSWER-REFERENCES-BEGIN
\section*{复习资料与使用说明}
本卷依据行列式章节的常见知识点重新编写，不对应某一教材的原题编号。
命题范围参考了以下公开课程的主题安排；各题计算与证明在本卷中独立给出。
\begin{itemize}
\item MIT OpenCourseWare，18.06SC：
\href{https://ocw.mit.edu/courses/18-06sc-linear-algebra-fall-2011/}{Linear Algebra}，
其中行列式相关单元涵盖性质、代数余子式、克拉默法则等内容。
\item MIT OpenCourseWare，18.013A：
\href{https://www.ocw.mit.edu/ans7870/18/18.013a/textbook/HTML/chapter32/section04.html}{32.4 More on Determinants}，
可用于复习行列式的线性性质、乘法公式和代数余子式。
\end{itemize}
上述网页核对日期为 2026 年 9 月 12 日。
% ANSWER-REFERENCES-END
\end{document}
